5.6.3.2 Main Effects

-----Section 1: DEX Scatter Plot-----

<label the plot as DEX Scatter Plot>

Conclusions:

1. Wide Spread: The most noteworthy aspect of the
   DEX scatter Plot is the wide range of the
   response: the distance ranges nearly 10 feet:
   from near-0 to beyond 125 inches.  Clearly,
   therefore, with such a range we have an
   interesting system: either we have a system in
   which the factors and their settings have no
   effect but is extremely noisy (unlikely), or we
   have a system in which the factors and settings
   have a deterministic effect which simply needs
   to be uncovered (more likely).

2. Bimodality: Also noteworthy is the bifurcation
   of the data into 2 regions: below 50 inches and
   above 80 inches, with a large 30+ inch gap in
   the middle.  Such bimodality is often caused by
   the presence of discrete (non-continuous)
   factors in the experiment, which would point an
   initial finger at factor 1 (band height) and
   factor 3 (number of bands) as the possible
   culprits.  In any event, the bimodality is
   glaring and so a side-goal of the analysis is
   to uncover the reason for its existence.

3. Settings for Y = 60: The gap from roughly 50 to
   80 is of particular concern because one of the
   stated engineering goals is to determine
   settings for the 5 factors that would yield a
   response of 60 inches, and there is no setting
   in any of the 20 runs that yielded a response
   anywhere close to Y = 60.  The implication is,
   therefore, that the best setting must come from
   a derived model and hence additional "pressure"
   is being put on the model that it have good
   interpolatory/extrapolatory properties.

4. Weak Main Factors: Note from the scatter plot
   that none of the 5 factors is (for one setting)
   consistently below and then (for the other
   setting) consistently above the gap.  This
   implies that there is no dominant factor and
   hence the resulting model will be more
   complicated and less parsimonious (that is will
   need many terms).

-----Section 2: Box Plot-----

An alternative is to generate a box plot, which
reduces all of 8 data points at a given factor
and setting into (in the simplest case) 5
summary statistics: the min, the 25% point, the
50% point (the median), the 75% point, and the
maximum.  Because of such simplicity, the box plot simetimes manifests
the differences in response locations (that is, the factor effects)
more clearly that the DEX scatter plot.

Conclusions:

1. Low Settings: Note how certain settings tend to
   yield low responses; in particular, smaller
   distances (<= 50 inches) tend to come from

          (X1,X2,X3,X4,X5) = (-,+,-,-,-).

   Further, these settings yield responses which
   also have low variability.  One can conclude,
   therefore, the the best setting to yield Y = 30
   (one of the experiment objectives) will
   probably be in the vicinity of (-,+,-,-,-).

2. High Settings: Note that the opposite settings:

          (X1,X2,X3,X4,X5) = (+,-,+,+,+)

   yield responses which are high location and
   high variability (that is, the distances are
   much longer, but the responses are also much
   noisier).  This implies that the Y = 90
   settings we develop will be much noisier (and
   hence of much poorer quality) that the Y = 30
   settings.

3. Variance-Stabilizing Transformations:
   Furthermore, this combination of low location &
   low variability with high location & high
   variability is characteristic of a system which
   violates the constant variability assumption
   implicit in subsequent least-squares model
   fitting, and hence provides the first
   suggestion that a variance-stabilizing
   transformation (e.g., logarithmic) should be
   applied to the data prior to further analysis.
   This transformation issue will be be revisited
   later.

-----Section 3: DEX Mean Plot-----

One of the experimental goals is to identify the most important
factors.  Specifically we would like to produce a ranked list
of such factors, ranked most important to least important.
Such "factors" can include main effects, but can in more generality
include interactions.  The DEX Mean Plot (also known as the "Main
Effects Plot) is an effective tool for producing such a ranked list
for main effect only (no interactions).

Note that the DEX scatter plot and the box plots had all the data
including center points.  For the DEX mean plot we produce 2 plots--
one with the center points and one without.  As is common in the
analysis of 2-level designs, we choose for subsequent analyses
to omit the center points and to use just the vertices of the 2**(5-1).
We hold the center points aside for analysis and for modeling and
we will use them at a later time for model validation/confirmation.

In a DEX Mean Plot, the vertical axis is the mean response at a setting
(typically this will be the mean of 8 values).
One virtue of orthogonla designs (such as the 2**(5-1) utilized here)
is that the least squares estimates
of the factor effects is identically the difference of the means for
each setting; hence, large differences (steep slopes) that manifest themselves
in a DEX Mean Plot implies large estimated factor effects.
The DEX Mean PLot hence provides an easy "eyeball" method for arriving
at an initial ranked list of (main) factors.

We produce 2 dex mean plots, one with the center points included
and one with them excluded.
.
Conclusions:

1. Non-Monotonicity: Note from the DEX Mean Plot with center points
   that for factors X3, X4, and X5, the center point has a higher
   mean response than at the -1 and +1 settings.
   The implication of this is that

      1. Curvature: The underlying response surface
         may have curvature and so a more complicated
         model may be needed.

      2. Inadequate model: Any fitted model based
         on the extreme points only may possibly
         fit poorly at the center point and
         therefore interpolation may be invalid.

      3. Biased effect estimates: The least
         squares estimates for the effects have an
         implicit assumption of response surface
         monotonicity (no response extrema between
         the extreme settings -1 and +1).  When
         this is violated, then the values of the
         effect estimates are extremely dependent
         on the settings, and may not
         monotonically increase/decrease as the
         settings are widened.  Worse, we may end
         up with effect estimates that are
         near-zero, when in fact the factor has a
         large effect on the response.

2. Ranked List of Factors: From the DEX Mean Plot with no center points,
   we may deduce the ranked list of important factors as

      1. X4 (arm length)      (effect = about 40 inches)
      2. X3 (number of bands) (effect = about 35 inches)
      3. X1 (band height)     (effect = about 30 inches)
      4. X5 (start point)     (effect = about 25 inches)
      5. X2 (stop angle)      (effect = about 20 inches)

3. Lack of Dominant Factors: From the DEX Mean
   Plot, we note that the factor effects drop off
   in a slow, near-uniform fashion--this is not
   typical.  It is much more common in engineering
   experiments to have 1 or 2 factors dominate and
   the remaining factor effects being 2 to 10
   times smaller--that was not the case here.
   There are 5 possible reasons for this
   non-dominance:

      1. Engineering reality
      2. Interactions & Confounding
      3. Other Factors
      4. Noisiness (high variance)
      5. Non-constant variance

   These 5 possibilities are now discussed point-by-point.

      1. Engineering Reality: The estimated effects
         may in fact be accurate, and this may be a
         factual reflection of the engineering
         cause-and-effect.  Again, this is not
         usually seen in engineering experiments, but
         is a possibility.

      2. Interactions & Confounding: Interactions may
         be important and hence (in fractional
         factorial designs) may contaminate main
         effects.  However, for the 2**(5-1)
         fractional factorial design that we have
         run, main effects are not confounded with
         any 2-term or 3-term interactions (only
         4-term interactions and higher).  For many
         physical systems, the assumption that 4-term
         interactions and higher are nil (and may be
         discarded) is invariably safe, and hence for
         this experiment, contamination via
         confounding is adjudged not to be the cause.

      3. Other factors: Other factors (besides the 5
         we have controlled) may in fact be driving
         the system.  Unfortunately, these other
         factors may be known or unknown, recorded or
         not recorded.  If there are other factors
         driving the system, then this is a
         near-fatal experimental deficiency, and their
         effect will be smeared over (that is, will
         bias) the main effects in an
         unpredictable way.  For this experiment,
         this possibility must still be entertained.
         In the larger issue, this is detected by the
         replicated (pseudo-) center points during
         the experiment, and the post-experiment
         confirmation runs (replicated) at the best
         settings (to yield Y = 30, 60, and 90).  For
         the replicated pseudo-center points, the
         spread is reasonably close (see below) and
         so there appears no drifiting caused by
         other factors.  Hence for this data, "other
         factors" is probably not the cause.

      4. Noisy Experiment: Some experiments are
         intrinsically noisy (for example, fracture and
         reliability data).  If the replication standard
         deviation (the intrinsic variability the data
         under fixed conditions) is high, then that
         may translate into large and non-discriminating
         effect estimates.  Fortunately, in our
         experiment design, we designed in some
         replication (the 2 pseudo-center points (each) at
         X3 = -1 and X3 = +1) so as estimate how much
         intrinsic replication error the catapult experiment
         has.  For the 2 pseudo-center points at
         (0,0,-1,0,0), the data is Y = 37.5 and 45, and
         hence the replication standard deviation s1 =
         5.303; for the 2 pseudo-center points at
         (0,0,1,0,0), the data is Y = 84.5 and 99 and so
         s2 = 9.175.  Combining these 2 estimates for sigma,
         we form a pooled standard deviation of s-pooled
         = 8.162.  Since 8.162 is less than 10% of the
         total spread of the data (8 to 126.5), "noisy
         experiment" is probably not the cause of the large
         effect estimates.

      5. Non-constant variance: This is the case where
         the data is partially noisy-- precise in
         certain subregions but noisy in other
         subregions.  In particular, for the case
         where the response variable Y is in a
         proportionately noisy (low response has
         low random error but large response has
         large random error), as opposed to absolutely
         noisy (all responses have same random error)
         then the setting of each factor which tend to
         yield large responses (distances) will be
         noisy, and hence tend to produce longer
         distances, larger means, and hence larger
         effect estimates.  For this experiment,
         non-constant variation is a plausible and
         the most likely explanation for the
         non-discriminating effect estimates.
         This discussion point will be picked up
         at a later time.

4. Settings to Maximize the Response: One of the
   engineering objectives of this experiment was
   to find the settings of the 5 factors that
   would yield Y = 30, 60, and 90.  If we
   temporarily set aside this stated objective,
   and replace it with the objective: "Find the
   settings which maximize the response" (which is
   common in many other engineering experiments),
   then the DEX Mean Plot assists in answering
   that question.  By noting for each factor as to
   whether the maximum mean is attained by -1 or
   +1, we may deduce from the plot that (on the
   average) the maximal response is attained by

          (X1,X2,X3,X4,X5) = (+1,-1,+1,+1,+1)

   Scanning the raw data table from section 5.6.3.1, we note
   that (+,-,+,+,+) happened by chance to not to be one of
   the design points (since this was a fractional factorial
   design).  However, the design point which yielded the
   observed maximum response (Y = 126.5) was (+,-,+,+,-).
   Note that these 2 sets of settings are identical for
   factors X1 to X4 and differ only in factor X5.  Based on
   the DEX Mean Plot, this suggests that (+,-,+,+,+) may in
   fact yield a larger response than the observed
   (+,-,+,+,-).  This may be confirmed (or denied) by
   follow-up confirmation runs at (+,-,+,+,+).

5. Settings to Yield Y = 30, 60, 90: As for the
   stated experimental objective of determining
   best settings to yield Y = 30, 60, and 90, we
   may deduce that the Y = 90 settings will be in
   the "vicinity" of the maximum setting:
   (+,-,+,+,+); the Y = 30 settings will be in the
   "vicinity" of the converse: (-,+,-,-,-); and the Y =
   60 settings are still unknown at this point.

-----Section 4: DEX Median Plot-----

Conclusions:

The DEX Median Plot is not usually done in the
recommended list of analysis techniques for
2-level orthogonal designs.  We employ it here
because of the aforementioned noting of
non-constant variation and how that ties in with
skewness.  We are specifically generating this
plot so as to compare its conclusions with the
conclusions of the DEX Mean Plot.  (Note, however,
for formal effect estimation (see section
5.6.3.6), we still employ the usual least squares
estimates which are the difference of the means.)
Definitionally, each plot point itself in the DEX
Median Plot is the median of all (8) points at the
given setting (-1 or +1) of the given factor.

Conclusions:

1. Skewness and Transformations: The DEX median
   plot indicates that the data is skewed.  Note
   from the DEX Mean Plot that the grand mean (of
   all 16 points) is approximately 55, whereas
   from the DEX Median Plot the grand median is
   drastically smaller: approximately 36.  Further
   the location for the higher points has dropped
   from around 70 (mean) to 60 (median), and the
   location of the lower points has dropped from
   around 40 (mean) to around 35 (median).  All of
   these are consistent with skewed data.  Skewed
   data is not fundamentally wrong to analyze; on
   the other hand, the usual analysis tools/tests
   are more valid with normal (or at least
   symmetric) data and so this suggests a
   transformation to normality (or at least
   symmetry).  Again, a logarithmic transformation
   is commonly useful to achieve such.

2. Ranked list of factors: Note that the ranked
   list of factors from the DEX Mean Plot was X4, X3,
   X1, X5, and X2, whereas the ranked list of
   factors from the DEX Median Plot is X4, X3, X5,
   X1, and X2.  The consistency and robustness of
   the X4 and X3 ranking across both plots
   provides confidence in the importance of
   these 2 factors relative to the remaining 3
   factors.

3. Best Settings for Maximum: Although this is not
   a stated objective of the experiment, note that
   the settings to yield the maximum (on the
   "median") is the same (+,-,+,+,+) as the
   settings to yield the maximum from the DEX Mean
   Plot.

-----Section 5: DEX Standard Deviation Plot-----

As with the DEX Median Plot, the DEX Standard
Deviation Plot is not usually generated (or
needed).  We employ it here to illustrate how it
may be usefully interpreted in certain anomalous
cases.  Each plot point itself is the simple
standard deviation of all (8) points at the given
setting (-1 or +1) of the given factor.

Conclusions:

1. Ranked List of "Noisy" Factors:
   Factors X3 (number of bands), X4 (arm length),
   and X1 (band height) all significantly differ
   in noisiness (by a factor of around 2 in the
   standard deviation) between the -1 and +1
   settings.  The +1 settings are about twice as
   noisy as th -1 setting for all 3 of these
   factors.  Factors X2 (start angle) and X5 (stop
   angle) have near-co variation across their 2
   settings.

2. Non-Constant Variation: It is useful to interpret
   the DEX standard deviation plot in conjunction
   with the DEX mean plot.  Note that the settings
   of the 5 factors which minimize the variability
   (from the DEX standard deviation plot) is
   (-,+,-,-,-).  Note that this is identical to
   the settings of the 5 factors which minimize
   the mean (from the DEX mean plot).  Similarly,
   large variability is from (+,-,+,+,+) which is
   identical to large location (+,-,+,+,+).
   Thus we a reinforcement of what was seen
   before, namely that the variation is increasing
   with the location--this is referred to as
   non-constant variation and its existence
   negates the usual coefficient estimation output
   from future least squares fitting.

-----Comparison of Plots-----

-----5.6.3.3 Interaction Effects-----

In addition to the main effects, it is also
important to check for interaction effects.  The
DEX mean interaction plot is an effective tool for
viewing akk 2-term interaction effects.  The
diagonal plots in the matrix are mean plots of each
main factor.  The off-diagonal plots are mean plots
versus each of the C(k,2) = C(5,2) = 10 2-term
interactions.  All plots have the same vertical axis
limits and so all are comparable.  Like usual, step
lines imply a strong effects while flat lines imply no
effects.  The least squares estimate of the effect is
given via annotation within each plot.  The matrix is
scanned by searching for steep lines (important
factors) to flat lines (unimportant factors).

Conclusions:

   1. Ranked List of Factors:
      1. X4  (arm length)      (effect =  40.3 inches)
      2. X3  (number of bands) (effect =  35.9 inches)
      3. X1  (band height)     (effect =  27.0 inches)
      4. X5  (start point)     (effect =  24.0 inches)
      5. X2  (stop angle)      (effect = -22.2 inches)
      6. X3*X4 interaction     (effect =  15.2 inches)
      7. X1*X4 interaction     (effect =   9.4 inches)
      8. X1*X3 interaction     (effect =   9.3 inches)
      etc.

   2. Interactions:  The most important interaction is X3*X4
      (number of bands)*(arm length).  The estimated interaction
      is 15.2 units which is more than 10% of the total spread
      (8 to 126.5) of the data, and hence is relatively large from
      a simple numeric point of view.  The 15.2 is to be interpreted
      as follows: this states that the X3 effect changes by 15.2
      units depending on whether X4 is -1 or +1.  The X3 effect of
      35.9 is a single number integrated over all of the data.
      Since the X3*x4 effect is 15.2 then this implies that the
      X3 effect for X4 = -1 equals 35.9 - 15.2/2 = 28.3 and
      the X3 effect for X4 = +1 will equal 35.9 + 15.2/2 = 43.5.

-----Block Plots-----

<as is>
Plots 1 and 2 address the question: Is factor X1 important?
The plot character (1 and 2) represent the 2 settings of factor 1.
The horizontal axis of plot 1 is the 4 combinations of X1 and X2.
The horizontal axis of plot 2 is the 4 combinations of X4 and X5.
Plots 3 and 4 ask: Is factor X2 important?
Plots 5 and 6 ask: Is factor X3 important?  etc.

Conclusions:

1. X1 is statistically significant (2 is larger than 1
   in 7 out of 8 cases--the probability of that (or worse)
   happening by chance is (8 + 1) / 256 = 3.5%.)
2. From plots 3 and 4, factor X2 has 2 > 1 in 3 out of
   8 cases--not statistically significant.
3. From plots 5 and 6, Factor X3 has 2 > 1 in 8 out of 8
   cases--statistically significant at the 1% level.
4. From plots 7 and 8, Factor X4 has 2 > 1 in 8 out of 8
   cases--also statistically significant.
5. From plots 9 and 10, Factor X5 has 2 > 1 in 4 out of 4
   cases--not statistically significant.
6. From plots 5 and 6, the height of the vertical bar
   (the local factor 3 estimated effect) is large
   when X1*X2 = -1 and small when X1*X2 = +1 thereby
   implying an X3*X1*X2 interaction.
7. etc.

   The block plot thus reinforces the conclusions from the
   Interaction effects matrix.

-----5.6.3.5: Model Selection, Fitting, & Validation-----

As with the Eddy Current Probe Sensitivity Full
Factorial Case Study, the criterion for fractional
factorial designs as to the ordering of terms to be
added in the model is the same, namely, add terms in
the order of the ranked list of important factors.
Hence our model will include X4, followed by X3,
followed by X1, etc.

When does one stop adding terms.  The guidelines are
the same as with full factorial disigns; namely,

1. Halfnormal Probability Plot of |effects|: add the
   terms which stand out in a halfnormal probability
   plot of the |effects|.

2. Minimum Engineering Signficant Difference: add those
   terms which are bigger than the engineer's
   prespecified minimum engineering significant
   difference (e.g., 10% of the total range of the data
   = 10% (126.5 - 8) = 12 inches = 1 foot.  (In fact,
   no such cutoff value was pre-specified in this
   experiment and so this criterion cannot be applied.)

3. Minimum Engineering Residual Standard Deviation: Add
   terms until the residual standard deviation gets
   smaller that the engineer's pre-specified value for
   how good he/she wants the model to be; that is, how
   small does he/she want the fitted model to be (e.g.,
   5% of the total range = 5% (126.5 - 8) = 6 inches).
   In fact, no such apriori value was pre-set and so
   this criterion cannot be applied.

4. Replication Standard Deviation: Add terms until the
   residual standard deviation is smaller than the
   replicated standard deviation.  Since this
   experiment has by design built-in replication, this
   criterion may be applied.  The replication standard
   deviations were gotten from the 2 pseudo-center
   points and had the values s1 = 5.6 and s2 = 9.x with
   a pooled value of 8.x.  The logic of this is that we
   can demand of our model to fit with more at least
   the precision of replicated points, but no further
   (else we would be fitting noise).  Thus keep adding
   terms to the model until the residual standard
   deviation drops below 8.x.

5. Normal Probability Plot of Residuals: Keep adding
   terms (and computing residuals) until the normal
   probability plot of the residuals is "sufficiently"
   linear.

The above 5 criterion may not all agree in terms of
what terms to include in the model.  Criterion 4 is an
absolute lower bound and so adding terms beyond that is
senseless.  Further, since this is essentially a
modeling exercise, then criterion 3 (how small do you
want the residual standard deviation to be?) is the
definitive criterion.  In the absense of an answer to
that question, then criterion 1 along with criterion 5
is often used in combination.

For the experiment at hand, we use criteria 1, 4, and 5.

The following plots show the normal probability plot of
the effects and the halfnormal probability plot of the
absolute value of the effects (both without and with
the coded effect tags).  The halfnormal probability
plot is more informative than the normal probability
plot since the halfnormal plot does not change
depending on how the original f settings were coded
(for example, for Factor 1, do we define 2.25 and 4.75
to be coded as -1 and +1 or conversely as +1 and -1?
Depending on what we do, this will change the sign of
the effect and hence change the appearance of the
normal probability plot--but the halfnormal plot will
be unchanged since it only focuses on magnitude).

<3 plots>

Conclusions:

1. The plots are not conclusive.  The normal
   probability plot is roughly linear which implies no
   factors are important.  The halfnormal probability
   plot is also roughly linear with the same
   implication.  In practice, such plots usually have a
   linear sub-portion (consisting of factors we leave
   out of the model) and then a displaced set of
   factors (which we include in the model).  Here the
   situation is not clear.  From the halfnormal
   probability plot, we can extract the upper 6 factors
   as (slightly) different from the lower linear 9
   factors, but little justification exists for this
   and less justification exists adding more terms
   beyond the 6. As from our ranked list, the 6 terms
   are X4, X3, X1, X5, X2, and X3*X4.

   Carrying out the usual least squares fit for this
   model, we obtain the following fitted model
   (prediction equation):

      Y = 55.30 + 0.5 (40.28 X4 + 35.91 X3 +
                       26.97 X1 + 24.09 X5 +
                      -22.16 X2 + 15.22 X3*X4)

   Regarding adequacy of this model, 2 approaches are
   carried out:

      1. Quantitatively: compute the residual standard
                         deviation and compare it to the
                         replication deviation (step 4 above);

      2. Graphically:    generate a normal probability plot
                         of the residuals so as to achieve
                         linearity (step 5 above).

   Quantitatively, the residual standard deviation for
   this model is 12.48.  This is 50% larger than the
   replication standard deviation of 8.xx and so this
   suggests that the model is inadequate (and so
   additional terms must be added).  Graphically, the
   normal probability plot of the residuals is
   presented below.

      <normal prob plot of residuals>

   This plot by itself suggests that the model is
   adequate, but the normal probability plot is a
   secondary tool relative to the computed residual
   standard deviation.  The bottom line is that the
   given model has a "typical residual" of the order of
   12 inches and so the usual +- 2 standard deviations
   will yield a rough 95% error of prediction of 24
   inches.  This is probably too large in practice and
   so this suggests that more terms need to be added to
   drive the error of prediction down.

   How many more terms need to be added?  To determine
   this we need the following table which provides
   effect estimates and residual standard deviations
   for cumulatively more complicated models.  One of
   the saving graces of orthogonal designs (such as the
   2**(5-1) fractional factorial design used here) is
   that the effect estimates do NOT change as we add
   additional terms.  The following table thus serves
   as a two-fold purpose of giving the ranked list of
   factors and also giving the goodness of fit of the
   correspoinding cumulative fitted models.

      <yates table here>

Conclusions:

As seen above, the model with a constant + 6 terms has
a residual standard deviation of 12.48.  Adding 2
additional terms (X1*X4 and X1*x3) brings the residual
standard deviation down to 10.02.  It would take 3 more
terms beyond that (a constant + 9 terms) to drive the
residual standard deviation below the replication
standard deviation of 8.16.  The normal probability
plot of the residuals of this model (not shown) would
be acceptable (that is, linear) but the problem we have
is 4-fold:

   1. Have we used the wrong form?  If the form (linear
      + interaction terms) is wrong, then no amount of
      additional terms will improve things.  For
      example, trying to fit an asymptotic exponential
      function via a polynomial.

   2. Is the model too complicated?  Such a model
      (constant + 9 terms) is always suspicious.  This
      suggests that perhaps another functional
      approach/form should have been used (e.g,
      transforming before fitting).

   3. Were the underlying assumptions for least squares
      fitting adhered to when estimating the
      coefficients of this model?  The answer to this
      is no--because of the non-constant variance
      problem (low responses have low variation, but
      high responses have high variation).  In such
      case, the regression coefficients of the model
      are incorrect.  The corrective action for this is
      twofold:

         1. weighted least squares;
         2. transformations of the response.

      This latter point will be developed and pursued
      in the next section.

   4. Does this model have good interpolatory
      properties?  Recall that the purpose of this
      model is not as an end in itself, but to assist
      in deriving good settings for Y = 30, 60, and 90.
      In this context, we ask as to how well does this
      model do at the pseudo-center points (such
      interpolatory testing is why these points were
      included in the design in the first place).

         At pseudo-center point 1: (0,0,-1,0,0), Y = 37.5 and 45.0,
            with a mean of 41.25.
            The fitted model yields a predicted value of
            Y = 55.30 + 0.5 (26.97*(-1)) = 41.82
            which is excellent compared to the mean of 41.25

         At pseudo-center point 2: (0,0,+1,0,0), Y = 84.5 and 99,
            with a mean of 91.75.
            The fitted model yields a predicted value of
            Y = 55.30 + 0.5 (26.97*(+1)) = 68.79
            which is terrible compared to the mean of 91.75

      We thus conclude that this model cannot be freely
      and universally used for interpolation--in
      particular for the X3 = +1 (that is, the number
      of rubber bands = 2) case.  We are thus led to
      the prospect of a model whose adequacy is
      conditional--based on the setting of a discrete
      variable (X3).  As a matter of experience, this
      is a quite common occurrance when discrete
      variables are involved in the analysis.  When
      this occurs (that such a factor is important,
      discrete, and the 2 levels of such a factor
      behave drastically different), then the
      suggestive corrective action is to split the
      analysis (and carry out 2 parallel analyses (and
      models, and settings)--one based on X3 = -1 and
      one based on X3 = +1.

      Four added benefits of such an approach are

         1. Simplified Sub-models: the resulting 2
            sub-models are each frequently quite simple
            (since all of the interaction terms
            involving X3 (in this case) disappear with
            the subsetting;

         2. Improved Predicted values: the prediction
            equations yield better (more accurate)
            predicted values;

         3. Improved best settings: This follows from
            the above.  This line of analysis will be
            pursued later.

         4. Improved Insight: Separating out a dominant
            discrete terms encourages a separation of
            engineering cause-effect into 2 separate
            camps, both devoid of complicated
            interactions;

Such subset fitting yields the following:

   For X3 = -1 data only (n = 8)

        Y =
